%----------------------------------------------------------------------------------------
% VIERKANTSWORTELS EN VERGELIJKINGEN. Alle opgaven zijn eigen werk.
%
% Opbouw: Start, leerstof (theorie, werkwijze of uitleg, voorbeeld, korte oefening), en op een
% nieuwe pagina de Opdrachten (door elkaar, stijgende moeilijkheid), Koppel en Test jezelf.
% Labels van dit hoofdstuk beginnen met h2-.
%----------------------------------------------------------------------------------------

\mychapter{2}{Vierkantswortels en vergelijkingen}

\lkr{Leerplan B+S': LPD 31 (tweedegraadsvergelijkingen met reële coëfficiënten oplossen in $\C$), met de wenk bij LPD 30 (vierkantswortels van een complex getal in cartesische vorm). De vergelijkingen van hogere graad in sectie 3 lossen leerlingen alleen op door af te zonderen en te ontbinden, met de technieken uit de 2de graad. Het wortelteken staat nooit bij een negatief getal.}

In hoofdstuk 1 bleek dat $2 + 3i$ een oplossing is van $x^2 - 4x + 13 = 0$, een vergelijking zonder reële oplossingen. Dat was geen toeval. In dit hoofdstuk zie je dat elke tweedegraadsvergelijking met reële coëfficiënten oplossingen heeft in $\C$, en hoe je ze vindt: met dezelfde discriminant als vroeger. De oplossingen komen in paren van toegevoegde getallen. Op het einde zoek je ook de vierkantswortels van een complex getal zoals $5 + 12i$.

\begin{oefeningen}[Start]
\oef{h2-start-1}

\oef{h2-verken}
\end{oefeningen}

\lkr{Laat oefening \ref{oef:h2-verken} maken vóór de uitleg. In vraag c lossen leerlingen een tweedegraadsvergelijking op zonder formule, door een kwadraat af te splitsen: dat is het bewijs van de stelling in sectie 2, met getallen. Vraag d bereidt som en product voor.}

%========================================================================================
\section{Vierkantswortels van een negatief getal}

\subsection{Een product is nul}

Bij het oplossen van vergelijkingen gebruik je in $\R$ voortdurend dat een product alleen 0 is als een van de factoren 0 is. Zo los je $(x - 2)(x + 5) = 0$ op zonder haakjes uit te werken. In dit hoofdstuk heb je dat ook in $\C$ nodig. Gelukkig blijft het gelden, dankzij het omgekeerde uit hoofdstuk 1.

\begin{theorie}
\begin{stelling}{Een product is nul}
Voor complexe getallen $z$ en $w$ is \quad $\formule{z \cdot w = 0 \iff z = 0 \text{ of } w = 0}$
\end{stelling}

\begin{bewijs}
Is $z = 0$ of $w = 0$, dan is het product 0, want $0 \cdot (c + di) = 0 + 0i = 0$. Stel omgekeerd $zw = 0$. Is $z = 0$, dan zijn we klaar. Is $z \neq 0$, dan heeft $z$ een omgekeerde $\frac{1}{z}$, en vermenigvuldigen we beide leden daarmee:
\bewijsregel{w = \frac{1}{z} \cdot z \cdot w = \frac{1}{z} \cdot 0 = 0 .}
\end{bewijs}
\end{theorie}

Een vergelijking die als product van factoren gelijk aan 0 staat, los je dus op door elke factor apart gelijk aan 0 te stellen, zoals in $\R$. Het verschil is dat de factoren nu complexe getallen mogen bevatten.

\voorbeeld{}{Los op in $\C$: $(2z - i)(iz + 3) = 0$.\label{vb:h2-product0}}{%
\uitleg{Het linkerlid is een product en het rechterlid is 0. We stellen elke factor gelijk aan 0, en lossen twee eerstegraadsvergelijkingen op zoals in hoofdstuk 1.}
\begin{stappen}
2z - i = 0 \quad \text{of} \quad iz + 3 = 0 & Een product is 0 als een factor 0 is. \\[1ex]
z = \frac{i}{2} \quad \text{of} \quad z = \frac{-3}{i} = \frac{-3 \cdot (-i)}{i \cdot (-i)} = 3i & Deel door 2, en door $i$ met de toegevoegde $-i$.
\end{stappen}
Controle van de tweede oplossing: $i \cdot 3i + 3 = 3i^2 + 3 = 0$.}

\begin{oefeningen}
\oef{h2-product0}
\end{oefeningen}

\subsection{Vierkantswortels van een negatief getal}

In $\R$ heeft een negatief getal geen vierkantswortels. In $\C$ wel: $3i$ en $-3i$ hebben allebei kwadraat $-9$. We breiden het begrip vierkantswortel daarom uit naar complexe getallen.

\begin{theorie}
Een complex getal $w$ is een \textbf{\textcolor{accent}{vierkantswortel}} van $z$ als \quad $\formule{w^2 = z}$

\begin{stelling}{Vierkantswortels van een negatief getal}
Voor een reëel getal $r > 0$ heeft $-r$ precies twee vierkantswortels: \quad $\formule{i\sqrt{r}}$ \ en \ $\formule{-i\sqrt{r}}$
\end{stelling}

\begin{bewijs}
Het zijn vierkantswortels, want $(\pm i\sqrt{r})^2 = i^2 \cdot \bigl(\sqrt{r}\bigr)^2 = -r$. Er zijn er niet meer: als $w^2 = -r$, dan is
\bewijsregel{w^2 - \bigl(i\sqrt{r}\bigr)^2 = 0, \quad \text{dus} \quad \bigl(w - i\sqrt{r}\bigr)\bigl(w + i\sqrt{r}\bigr) = 0 ,}
en een product is alleen 0 als een van de factoren 0 is, dus $w = i\sqrt{r}$ of $w = -i\sqrt{r}$.
\end{bewijs}
\end{theorie}

Elk negatief getal heeft dus twee vierkantswortels, net zoals een positief getal: twee tegengestelde getallen. Bij een positief getal liggen ze op de reële as, bij een negatief getal op de imaginaire as.

We schrijven de vierkantswortels als $\pm i\sqrt{5}$ en nooit als $\sqrt{-5}$. Het wortelteken bij een negatief getal leidt tot fouten. Wie $\sqrt{-4} \cdot \sqrt{-9} = \sqrt{36} = 6$ schrijft, vergist zich: met $2i$ en $3i$ is het product $6i^2 = -6$. De regel $\sqrt{a} \cdot \sqrt{b} = \sqrt{ab}$ geldt alleen voor $a, b \geqslant 0$. Je schrijft $i$ liefst vóór de wortel, zodat je niet leest dat $i$ onder het wortelteken staat.

\voorbeeld{}{Bepaal de vierkantswortels van $-12$.\label{vb:h2-wortel}}{%
\uitleg{Het getal is negatief, dus de vierkantswortels zijn $\pm i\sqrt{12}$. We vereenvoudigen de wortel.}
\begin{stappen}
\pm i\sqrt{12} = \pm i\sqrt{4 \cdot 3} = \pm 2i\sqrt3 & $\sqrt{12} = 2\sqrt3$. \\[1ex]
\bigl(2i\sqrt3\bigr)^2 = 4 \cdot i^2 \cdot 3 = -12 & Controle.
\end{stappen}}

\begin{oefeningen}
\oef{h2-wortel}
\end{oefeningen}

%========================================================================================
\section{Tweedegraadsvergelijkingen}

\subsection{De discriminant}

Een tweedegraadsvergelijking $az^2 + bz + c = 0$ los je op met de discriminant $D = b^2 - 4ac$. Was $D < 0$, dan was er in $\R$ geen oplossing, omdat je de vierkantswortel van $D$ nodig had. In $\C$ heeft $D$ wel vierkantswortels. In de Verken heb je dat al gebruikt: je splitste een kwadraat af, en zocht de vierkantswortels van een negatief getal. Het bewijs doet hetzelfde met letters.

\begin{theorie}
\begin{stelling}{Tweedegraadsvergelijking met reële coëfficiënten}
De vergelijking $az^2 + bz + c = 0$ met $a, b, c \in \R$, $a \neq 0$ en discriminant $D = b^2 - 4ac$ heeft in $\C$

\vspace{1mm}
als $D \geqslant 0$: de reële oplossingen \quad $\formule{z = \dfrac{-b \pm \sqrt{D}}{2a}}$

\vspace{1mm}
als $D < 0$: de twee toegevoegde oplossingen \quad $\formule{z = \dfrac{-b \pm i\sqrt{-D}}{2a}}$
\end{stelling}

\begin{bewijs}
Deel door $a$ en splits een kwadraat af, zoals in de Verken:
\bewijsregel{z^2 + \frac{b}{a}z + \frac{c}{a} = 0 \iff \Bigl(z + \frac{b}{2a}\Bigr)^2 - \frac{b^2}{4a^2} + \frac{c}{a} = 0 \iff \Bigl(z + \frac{b}{2a}\Bigr)^2 = \frac{D}{4a^2} .}
Neem een getal $\delta$ met $\delta^2 = D$: \ $\delta = \sqrt{D}$ als $D \geqslant 0$, en $\delta = i\sqrt{-D}$ als $D < 0$. Dan staat er
\bewijsregel{\Bigl(z + \frac{b}{2a}\Bigr)^2 = \Bigl(\frac{\delta}{2a}\Bigr)^2 , \quad \text{dus} \quad \Bigl(z + \frac{b}{2a} - \frac{\delta}{2a}\Bigr)\Bigl(z + \frac{b}{2a} + \frac{\delta}{2a}\Bigr) = 0 .}
Een product is 0 als een factor 0 is, dus $z = \frac{-b + \delta}{2a}$ of $z = \frac{-b - \delta}{2a}$. Andere oplossingen zijn er niet. Is $D < 0$, dan zijn dat de getallen $\frac{-b}{2a} \pm \frac{\sqrt{-D}}{2a}\,i$, elkaars toegevoegde.
\end{bewijs}
\end{theorie}

Je hoeft dus maar één formule te kennen: als $D < 0$ vervang je $\sqrt{D}$ door $i\sqrt{-D}$. In $\C$ heeft een tweedegraadsvergelijking met reële coëfficiënten altijd oplossingen: twee als $D \neq 0$, één als $D = 0$. Het bewijs toont ook dat er nooit meer zijn: wie in $\C$ zoekt, vindt bij $D \geqslant 0$ geen extra niet-reële oplossingen.

\voorbeeld{}{Los op in $\C$: $z^2 - 4z + 13 = 0$.\label{vb:h2-vkv}}{%
\uitleg{We berekenen de discriminant. Die is negatief, dus vervangen we $\sqrt{D}$ door $i\sqrt{-D}$.}
\begin{stappen}
D = (-4)^2 - 4 \cdot 1 \cdot 13 = 16 - 52 = -36 & $D < 0$: twee toegevoegde oplossingen. \\[1ex]
z = \frac{4 \pm i\sqrt{36}}{2} = \frac{4 \pm 6i}{2} & De vierkantswortels van $-36$ zijn $\pm 6i$. \\[3ex]
z = 2 + 3i \quad \text{of} \quad z = 2 - 3i & Deel reëel en imaginair deel door 2.
\end{stappen}}

\voorbeeld{}{Los op in $\C$: $2z^2 + 2z + 5 = 0$.\label{vb:h2-vkv-a}}{%
\uitleg{Nu is $a = 2$: let op de noemer $2a = 4$.}
\begin{stappen}
D = 2^2 - 4 \cdot 2 \cdot 5 = 4 - 40 = -36 & Weer $D < 0$. \\[1ex]
z = \frac{-2 \pm 6i}{4} = -\frac12 \pm \frac32 i & Deel elke term van de teller door 4.
\end{stappen}}

Ontbreekt de term in $z$, dan heb je de formule niet nodig: je zondert $z^2$ af, en zoekt de vierkantswortels van een negatief getal.

\voorbeeld{}{Los op in $\C$: $25z^2 + 16 = 0$.\label{vb:h2-vkv-b}}{%
\uitleg{Hier is $b = 0$. We zonderen $z^2$ af, zoals bij een vergelijking $x^2 = 4$ in $\R$.}
\begin{stappen}
z^2 = -\frac{16}{25} & Trek 16 af en deel door 25. \\[2ex]
z = \pm i\sqrt{\frac{16}{25}} = \pm\frac45 i & De vierkantswortels van een negatief getal.
\end{stappen}
Met de formule vind je hetzelfde: $D = -1600$ en $z = \frac{\pm 40i}{50} = \pm\frac45 i$.}

\begin{werkwijze}
Breng alles naar één lid en bereken $D = b^2 - 4ac$. Is $D < 0$, zoek dan de vierkantswortels $\pm i\sqrt{-D}$ en vul ze in in de formule, in de plaats van $\pm\sqrt{D}$. Schrijf de oplossingen in de vorm $a + bi$. Is $b = 0$, dan zonder je $z^2$ af.
\end{werkwijze}

\nujij{Los op in $\C$: $z^2 + 2z + 2 = 0$.}{$D = 4 - 8 = -4$, dus $z = \frac{-2 \pm 2i}{2} = -1 \pm i$.}

\naast{\begin{minipage}{10cm}
Waar zitten die oplossingen in een grafiek? De parabool met vergelijking $y = x^2 - 4x + 13$ snijdt de $x$-as niet: daarom zijn er geen reële oplossingen. De complexe oplossingen $2 \pm 3i$ zie je niet als snijpunten, maar hun reëel deel 2 is wel de $x$-coördinaat van de top. Dat is altijd zo, want het reëel deel van de oplossingen is $-\frac{b}{2a}$. In Koppel ontdek je wat het imaginair deel met de grafiek te maken heeft.
\end{minipage}}{\begin{tikzpicture}[x=.5cm,y=.2cm,font=\scriptsize]
  \draw[black!14,very thin,xstep=1,ystep=5] (-1,0) grid (5,20);
  \draw[->,grijs] (-1,0) -- (5.5,0) node[right,text=black] {$x$};
  \draw[->,grijs] (0,0) -- (0,21.5) node[above,text=black] {$y$};
  \foreach \t in {2,4}{\node[below,inner sep=1.5pt] at (\t,0) {$\t$};}
  \foreach \t in {5,10,15,20}{\node[left,inner sep=1.5pt] at (0,\t) {$\t$};}
  \draw[accent,thick] plot[domain=-.6:4.6,samples=60,smooth] (\x,{(\x)^2-4*(\x)+13});
  \fill[accent] (2,9) circle (2pt) node[below,text=black,inner sep=3pt] {top $(2, 9)$};
\end{tikzpicture}}

\begin{oefeningen}
\oef{h2-vkv}

\oef{h2-vkv2}
\end{oefeningen}

%========================================================================================
\section{Toegevoegde oplossingen}

\subsection{Rekenregels voor de toegevoegde}

De twee niet-reële oplossingen van een tweedegraadsvergelijking met reële coëfficiënten zijn altijd elkaars toegevoegde. Dat volgt uit de formule, maar het geldt ook voor vergelijkingen van hogere graad. Om dat te bewijzen, hebben we eerst rekenregels voor de toegevoegde nodig. We bekijken ze één voor één.

Wat gebeurt er met de toegevoegde als je twee getallen optelt? Neem je eerst de som en dan de toegevoegde, of eerst de toegevoegden en dan de som?

\begin{theorie}
\begin{stelling}{Toegevoegde van een som}
Voor complexe getallen $z$ en $w$ is \quad $\formule{\conj{z + w} = \conj{z} + \conj{w}}$
\end{stelling}

\begin{bewijs}
Neem $z = a + bi$ en $w = c + di$. Dan is $z + w = (a + c) + (b + d)i$, en dus
\bewijsregel{\conj{z + w} = (a + c) - (b + d)i = (a - bi) + (c - di) = \conj{z} + \conj{w} .}
\end{bewijs}
\end{theorie}

De volgorde maakt dus niet uit. Hetzelfde geldt voor een som van meer dan twee getallen: pas de regel een paar keer na elkaar toe. Bij een product is dat minder vanzelfsprekend, want daar vermengen het reëel en het imaginair deel zich.

\begin{theorie}
\begin{stelling}{Toegevoegde van een product}
Voor complexe getallen $z$ en $w$ is \quad $\formule{\conj{z \cdot w} = \conj{z} \cdot \conj{w}}$
\end{stelling}

\begin{bewijs}
Neem $z = a + bi$ en $w = c + di$. Voor het product is $zw = (ac - bd) + (ad + bc)i$, dus
\bewijsregel{\conj{z \cdot w} = (ac - bd) - (ad + bc)i .}
Anderzijds is
\bewijsregel{\conj{z} \cdot \conj{w} = (a - bi)(c - di) = ac - adi - bci + bdi^2 = (ac - bd) - (ad + bc)i .}
Beide uitkomsten zijn gelijk.
\end{bewijs}
\end{theorie}

Een reëel getal $r$ is zijn eigen toegevoegde: $\conj{r} = r$, want zijn imaginair deel is 0. Met de regel voor het product geeft dat $\conj{r \cdot z} = r \cdot \conj{z}$: een reële factor blijft staan.

\voorbeeld{}{Controleer de regel voor het product met $z = 3 + i$ en $w = 1 - 2i$.\label{vb:h2-conjproduct}}{%
\uitleg{We rekenen beide leden apart uit: eerst het product en dan de toegevoegde, daarna eerst de toegevoegden en dan het product.}
\begin{stappenbreed}
zw = 3 - 6i + i - 2i^2 = 5 - 5i & Het product. \\[1ex]
\conj{zw} = 5 + 5i & De toegevoegde ervan. \\[1ex]
\conj{z} \cdot \conj{w} = (3 - i)(1 + 2i) = 3 + 6i - i - 2i^2 = 5 + 5i & Hetzelfde getal.
\end{stappenbreed}}

Een macht is een product van gelijke factoren. Daarom volgt de regel voor een macht uit die voor een product.

\begin{theorie}
\begin{stelling}{Toegevoegde van een macht}
Voor een complex getal $z$ en een natuurlijk getal $n$ is \quad $\formule{\conj{z^n} = \conj{z}^{\,n}}$
\end{stelling}

\begin{bewijs}
Voor $n = 0$ staat er $\conj{1} = 1$, en voor $n = 1$ staat er $\conj{z} = \conj{z}$. Voor $n = 2$ geeft de regel voor het product
\bewijsregel{\conj{z^2} = \conj{z \cdot z} = \conj{z} \cdot \conj{z} = \conj{z}^{\,2} .}
Zo gaat het verder: geldt de regel voor een exponent $n$, dan is $\conj{z^{n+1}} = \conj{z^n \cdot z} = \conj{z^n} \cdot \conj{z} = \conj{z}^{\,n} \cdot \conj{z} = \conj{z}^{\,n+1}$, en geldt hij ook voor $n + 1$.
\end{bewijs}
\end{theorie}

Ken je dus een macht van een getal, dan ken je meteen dezelfde macht van zijn toegevoegde. Dat bespaart veel rekenwerk.

\voorbeeld{}{Gegeven: $(2 + i)^5 = -38 + 41i$. Bereken $(2 - i)^5$.\label{vb:h2-conjmacht}}{%
\uitleg{Het getal $2 - i$ is de toegevoegde van $2 + i$. We hoeven dus niets uit te werken: de toegevoegde van de macht is de macht van de toegevoegde.}
\begin{stappen}
(2 - i)^5 = \bigl(\conj{2 + i}\bigr)^5 = \conj{(2 + i)^5} & De regel voor een macht, van rechts naar links gelezen. \\[1ex]
= \conj{-38 + 41i} = -38 - 41i & Verander het teken van het imaginair deel.
\end{stappen}}

\begin{oefeningen}
\oef{h2-conj}
\end{oefeningen}

\subsection{Toegevoegde oplossingen}

Met de drie rekenregels kunnen we nu bewijzen dat de niet-reële oplossingen van een veeltermvergelijking met reële coëfficiënten altijd in toegevoegde paren voorkomen, welke graad de vergelijking ook heeft.

\begin{theorie}
\begin{stelling}{Toegevoegde oplossingen}
Is $z$ een oplossing van een veeltermvergelijking met reële coëfficiënten
\bewijsregel{\formule{a_n z^n + a_{n-1} z^{n-1} + \cdots + a_1 z + a_0 = 0} \qquad (a_0, a_1, \ldots, a_n \in \R),}
dan is ook $\formule{\conj{z}}$ een oplossing.
\end{stelling}

\begin{bewijs}
Neem van beide leden de toegevoegde. Rechts staat $\conj{0} = 0$. Links is de toegevoegde van de som de som van de toegevoegden, en voor elke term is
\bewijsregel{\conj{a_k z^k} = \conj{a_k} \cdot \conj{z^k} = a_k \cdot \conj{z}^{\,k} ,}
met de regels voor het product en voor een macht, en omdat $a_k$ reëel is. Samen geeft dat
\bewijsregel{a_n \conj{z}^{\,n} + a_{n-1} \conj{z}^{\,n-1} + \cdots + a_1 \conj{z} + a_0 = 0 ,}
dus $\conj{z}$ is ook een oplossing.
\end{bewijs}
\end{theorie}

Ken je dus één niet-reële oplossing, dan ken je er twee. In het vlak van Gauss liggen de oplossingen van zo'n vergelijking symmetrisch ten opzichte van de reële as.

\voorbeeld{}{Ga na dat $1 + i$ een oplossing is van $z^4 + 4 = 0$, en geef zonder te rekenen nog een oplossing.\label{vb:h2-toegevoegd1}}{%
\uitleg{We vullen in om te controleren. De coëfficiënten 1 en 4 zijn reëel, dus de toegevoegde is ook een oplossing.}
\begin{stappen}
(1 + i)^2 = 1 + 2i + i^2 = 2i & Eerst het kwadraat. \\[1ex]
(1 + i)^4 = (2i)^2 = 4i^2 = -4 & Dus $(1 + i)^4 + 4 = 0$. \\[1ex]
1 - i & De toegevoegde is ook een oplossing.
\end{stappen}
Controle: $(1 - i)^2 = -2i$ en $(-2i)^2 = -4$.}

De voorwaarde dat de coëfficiënten reëel zijn, is nodig. De vergelijking $z^2 - 3iz - 2 = 0$ heeft de niet-reële coëfficiënt $-3i$. Ze is te schrijven als $(z - i)(z - 2i) = 0$: werk de haakjes uit om dat na te gaan. De oplossingen $i$ en $2i$ zijn geen toegevoegden, want de toegevoegde van $i$ is $-i$. In opdracht \ref{oef:h2-o15} zie je nog zo'n vergelijking.

\subsection{Som en product van de oplossingen}

In de Verken was de som van de oplossingen $-3 + 2i$ en $-3 - 2i$ gelijk aan $-6$, en hun product $13$: de coëfficiënten van $z^2 + 6z + 13$, op het teken na. Dat ken je al bij reële oplossingen, en het blijft gelden in $\C$. We bewijzen het met de formule, met dezelfde $\delta$ als in het bewijs van de discriminant: $\delta^2 = D$.

\begin{theorie}
\begin{stelling}{Som van de oplossingen}
Heeft $az^2 + bz + c = 0$ de oplossingen $z_1$ en $z_2$, dan is \quad $\formule{z_1 + z_2 = -\dfrac{b}{a}}$
\end{stelling}

\begin{bewijs}
Volgens de formule is $z_1 = \frac{-b + \delta}{2a}$ en $z_2 = \frac{-b - \delta}{2a}$, dus
\bewijsregel{z_1 + z_2 = \frac{-b + \delta - b - \delta}{2a} = \frac{-2b}{2a} = -\frac{b}{a} .}
\end{bewijs}

\begin{stelling}{Product van de oplossingen}
Heeft $az^2 + bz + c = 0$ de oplossingen $z_1$ en $z_2$, dan is \quad $\formule{z_1 \cdot z_2 = \dfrac{c}{a}}$
\end{stelling}

\begin{bewijs}
Met het merkwaardig product $(x + y)(x - y) = x^2 - y^2$ en $\delta^2 = D = b^2 - 4ac$ is
\bewijsregel{z_1 z_2 = \frac{(-b + \delta)(-b - \delta)}{4a^2} = \frac{b^2 - \delta^2}{4a^2} = \frac{b^2 - b^2 + 4ac}{4a^2} = \frac{c}{a} .}
\end{bewijs}
\end{theorie}

Is $D = 0$, dan is er maar één oplossing: neem dan $z_1 = z_2$, en de formules gelden nog. Bij toegevoegde oplossingen $p \pm qi$ is de som $2p$ en het product $(p + qi)(p - qi) = p^2 + q^2$, de modulus in het kwadraat. Zo controleer je snel je oplossingen, of stel je een vergelijking op.

\voorbeeld{}{Bepaal de som en het product van de oplossingen van $2z^2 - 6z + 9 = 0$ zonder de vergelijking op te lossen. Controleer daarna door op te lossen.\label{vb:h2-somproduct}}{%
\uitleg{Met $a = 2$, $b = -6$ en $c = 9$ vullen we de formules in. Voor de controle lossen we op met de discriminant.}
\begin{stappen}
z_1 + z_2 = -\frac{-6}{2} = 3, \quad z_1 z_2 = \frac92 & Let op het minteken bij de som. \\[3ex]
D = 36 - 72 = -36, \quad z = \frac{6 \pm 6i}{4} = \frac32 \pm \frac32 i & De oplossingen. \\[3ex]
\frac32 + \frac32 = 3, \quad \Bigl(\frac32\Bigr)^2 + \Bigl(\frac32\Bigr)^2 = \frac94 + \frac94 = \frac92 & Som $2p$ en product $p^2 + q^2$: het klopt.
\end{stappen}}

\begin{oefeningen}
\oef{h2-somproduct}
\end{oefeningen}

Omgekeerd kun je uit één niet-reële oplossing een vergelijking opstellen: de andere oplossing is de toegevoegde, en met de som en het product ken je de coëfficiënten.

\voorbeeld{}{Stel een tweedegraadsvergelijking met reële coëfficiënten op die $1 - 4i$ als oplossing heeft.\label{vb:h2-opstellen}}{%
\uitleg{De coëfficiënten zijn reëel, dus de andere oplossing is de toegevoegde. Met de som en het product kennen we de vergelijking.}
\begin{stappen}
z_2 = 1 + 4i & De toegevoegde oplossing. \\[1ex]
z_1 + z_2 = 2, \quad z_1 z_2 = 1^2 + 4^2 = 17 & Som en product. \\[1ex]
z^2 - 2z + 17 = 0 & Met $a = 1$: $b = -2$ en $c = 17$.
\end{stappen}
Elk veelvoud, zoals $3z^2 - 6z + 51 = 0$, is ook goed.}

\nujij{Stel een vergelijking op met oplossing $5i$.}{De andere oplossing is $-5i$; som 0, product 25: $z^2 + 25 = 0$.}

\subsection{Ontbinden in factoren}

In $\R$ kun je $x^2 + 2x + 10$ niet ontbinden in factoren van de eerste graad, omdat de discriminant negatief is. In $\C$ lukt dat wel, met de oplossingen van de vergelijking.

\begin{theorie}
\begin{stelling}{Ontbinding in factoren}
Heeft $az^2 + bz + c = 0$ de oplossingen $z_1$ en $z_2$, dan is

\vspace{1mm}
\qquad $\formule{az^2 + bz + c = a(z - z_1)(z - z_2)}$
\end{stelling}

\begin{bewijs}
Werk het rechterlid uit, en gebruik de som en het product van de oplossingen:
\bewijsregel{a(z - z_1)(z - z_2) = a\bigl(z^2 - (z_1 + z_2)z + z_1 z_2\bigr) = a\Bigl(z^2 + \frac{b}{a}z + \frac{c}{a}\Bigr) = az^2 + bz + c .}
\end{bewijs}
\end{theorie}

In $\C$ kun je dus elke tweedegraadsveelterm met reële coëfficiënten ontbinden in twee factoren van de eerste graad. Vergeet de factor $a$ niet als $a \neq 1$.

\voorbeeld{}{Ontbind $z^2 + 2z + 10$ in factoren van de eerste graad.\label{vb:h2-ontbinden}}{%
\uitleg{In $\R$ lukt dat niet, want de discriminant is negatief. In $\C$ zoeken we de oplossingen van $z^2 + 2z + 10 = 0$, en gebruiken we de ontbinding uit het kader.}
\begin{stappen}
D = 4 - 40 = -36, \quad z = \frac{-2 \pm 6i}{2} = -1 \pm 3i & De oplossingen. \\[3ex]
z^2 + 2z + 10 = (z + 1 - 3i)(z + 1 + 3i) & $z - (-1 + 3i) = z + 1 - 3i$.
\end{stappen}
Controle: $(z + 1)^2 - (3i)^2 = z^2 + 2z + 1 + 9$.}

\begin{oefeningen}
\oef{h2-opstellen}
\end{oefeningen}

\subsection{Vergelijkingen van hogere graad}

Een vergelijking van hogere graad los je op zoals in $\R$: je brengt ze terug tot vergelijkingen van lagere graad, door een factor af te zonderen of te ontbinden. Het verschil is dat een factor als $z^2 + 4$ nu ook oplossingen geeft. We bekijken vier manieren om te ontbinden.

\voorbeeld{}{Los op in $\C$: $z^3 + z^2 + 4z + 4 = 0$.\label{vb:h2-hoger}}{%
\uitleg{We ontbinden door termen te groeperen: de eerste twee termen hebben $z^2$ gemeen, de laatste twee 4.}
\begin{stappen}
z^2(z + 1) + 4(z + 1) = 0 & Zonder in de groepen een factor af. \\[1ex]
(z + 1)(z^2 + 4) = 0 & De factor $z + 1$ komt twee keer voor. \\[1ex]
z = -1 \quad \text{of} \quad z^2 = -4 & Een product is 0 als een factor 0 is. \\[1ex]
z = -1 \quad \text{of} \quad z = 2i \quad \text{of} \quad z = -2i & Drie oplossingen, waarvan twee toegevoegd.
\end{stappen}}

Lukt groeperen niet, zoek dan een reële oplossing $r$ tussen de delers van de constante term. Dan kun je de veelterm delen door $z - r$, bijvoorbeeld met de regel van Horner.

\voorbeeld{}{Los op in $\C$: $z^3 - 3z^2 + 4z - 2 = 0$.\label{vb:h2-horner}}{%
\uitleg{We proberen de delers van de constante term $-2$. Vinden we een oplossing $r$, dan delen we door $z - r$, en lossen we het quotiënt op.}
\begin{stappen}
1 - 3 + 4 - 2 = 0 & Invullen van $z = 1$: het is een oplossing. \\[1ex]
z^3 - 3z^2 + 4z - 2 = (z - 1)(z^2 - 2z + 2) & Horner met 1 geeft de coëfficiënten $1$, $-2$ en $2$, en rest 0. \\[1ex]
D = 4 - 8 = -4, \quad z = \frac{2 \pm 2i}{2} = 1 \pm i & De tweedegraadsfactor. \\[3ex]
z = 1 \quad \text{of} \quad z = 1 + i \quad \text{of} \quad z = 1 - i &
\end{stappen}}

Komen alleen even machten van $z$ voor, dan stel je $t = z^2$, zoals bij de vergelijking in de Opwarmer. In $\R$ viel de negatieve waarde van $t$ weg; in $\C$ geeft ze twee oplossingen.

\voorbeeld{}{Los op in $\C$: $z^4 - 5z^2 - 36 = 0$.\label{vb:h2-bikwadratisch}}{%
\uitleg{Stel $t = z^2$. Dan staat er een tweedegraadsvergelijking in $t$. Voor elke oplossing $t$ zoeken we daarna de vierkantswortels.}
\begin{stappen}
t^2 - 5t - 36 = 0 & Met $t = z^2$. \\[1ex]
t = 9 \quad \text{of} \quad t = -4 & $D = 25 + 144 = 169$, dus $t = \frac{5 \pm 13}{2}$. \\[1ex]
z^2 = 9 \quad \text{of} \quad z^2 = -4 & Terug naar $z$. \\[1ex]
z = \pm 3 \quad \text{of} \quad z = \pm 2i & In $\R$ waren er twee oplossingen, in $\C$ vier.
\end{stappen}}

Ken je een niet-reële oplossing, dan ken je er meteen twee: de toegevoegde is ook een oplossing. Hun product met de twee factoren geeft een tweedegraadsfactor met reële coëfficiënten.

\voorbeeld{}{Het getal $2 + i$ is een oplossing van $z^3 - 3z^2 + z + 5 = 0$. Bepaal alle oplossingen.\label{vb:h2-toegevoegd}}{%
\uitleg{De coëfficiënten zijn reëel, dus ook $2 - i$ is een oplossing. De twee factoren $z - (2 + i)$ en $z - (2 - i)$ geven samen een reële tweedegraadsfactor, en de derde factor vinden we door de constante termen te vergelijken.}
\begin{stappenbreed}
z = 2 - i & De toegevoegde oplossing. \\[1ex]
(z - 2 - i)(z - 2 + i) = (z - 2)^2 - i^2 = z^2 - 4z + 5 & Merkwaardig product. \\[1ex]
z^3 - 3z^2 + z + 5 = (z^2 - 4z + 5)(z + c) & De andere factor is van de eerste graad. \\[1ex]
5c = 5, \quad c = 1 & Vergelijk de constante termen. \\[1ex]
z = 2 + i \quad \text{of} \quad z = 2 - i \quad \text{of} \quad z = -1 & De factor $z + 1$ geeft $z = -1$.
\end{stappenbreed}
Controle: $(z^2 - 4z + 5)(z + 1) = z^3 - 3z^2 + z + 5$.}

\begin{werkwijze}
Ontbind de veelterm: zonder een factor af, groepeer termen, stel $t = z^2$ als er alleen even machten staan, of deel door $z - r$ als je een oplossing $r$ kent. Ken je een niet-reële oplossing, gebruik dan ook de toegevoegde. Los daarna elke factor op, ook de tweedegraadsfactoren met $D < 0$.
\end{werkwijze}

\begin{oefeningen}
\oef{h2-hoger}

\oef{h2-hoger2}

\oef{h2-hoger3}
\end{oefeningen}

%========================================================================================
\section{Vierkantswortels van een complex getal}

Elk negatief getal heeft twee vierkantswortels in $\C$. Hoe zit het met een niet-reëel getal, zoals $5 + 12i$? In de Opwarmer heb je $(3 + 2i)^2 = 5 + 12i$ uitgerekend, dus $3 + 2i$ is een vierkantswortel van $5 + 12i$, en $-(3 + 2i)$ ook. Heeft elk complex getal zo'n vierkantswortel, en hoe vind je die zonder ze te raden?

\begin{theorie}
\begin{stelling}{Vierkantswortels van een complex getal}
Elk complex getal $z \neq 0$ heeft precies twee vierkantswortels $\formule{w}$ en $\formule{-w}$.
\end{stelling}

\begin{bewijs}
Neem $z = a + bi$ en $r = |z| = \sqrt{a^2 + b^2}$. Omdat $r^2 = a^2 + b^2 \geqslant a^2$, is $r \geqslant |a|$, en zijn $r + a$ en $r - a$ niet negatief. Stel daarom
\bewijsregel{x = \sqrt{\frac{r + a}{2}} \quad \text{en} \quad y = \pm\sqrt{\frac{r - a}{2}} ,}
met het plusteken als $b \geqslant 0$ en het minteken als $b < 0$. Dan is $x^2 - y^2 = \frac{r + a}{2} - \frac{r - a}{2} = a$, en
\bewijsregel{2xy = \pm 2\sqrt{\frac{(r + a)(r - a)}{4}} = \pm\sqrt{r^2 - a^2} = \pm\sqrt{b^2} = \pm|b| = b ,}
omdat het teken dat van $b$ is. Dus $w = x + yi$ is een vierkantswortel: $w^2 = (x^2 - y^2) + 2xyi = a + bi$. Ook $(-w)^2 = w^2 = z$. Zijn er meer? Is $v^2 = z = w^2$, dan is $(v - w)(v + w) = 0$, dus $v = w$ of $v = -w$. En $w \neq -w$, want $w \neq 0$.
\end{bewijs}
\end{theorie}

Elk complex getal verschillend van 0 heeft dus twee tegengestelde vierkantswortels, net als een positief of een negatief reëel getal. In hoofdstuk 4 vind je ze ook meetkundig.

Je hoeft de formules uit het bewijs niet van buiten te kennen. In de praktijk stel je $w = x + yi$ met $x, y \in \R$, en schrijf je de voorwaarde $w^2 = z$ uit:
\[ w^2 = (x^2 - y^2) + 2xyi . \]
Dat geeft twee vergelijkingen: één voor het reëel deel en één voor het imaginair deel. Een derde vergelijking, die het rekenwerk eenvoudig maakt, haal je uit de modulus. De modulus van $w^2$ is
\[ |w^2| = \sqrt{(x^2 - y^2)^2 + (2xy)^2} = \sqrt{x^4 + 2x^2y^2 + y^4} = \sqrt{(x^2 + y^2)^2} = x^2 + y^2 , \]
en uit $w^2 = z$ volgt dus $x^2 + y^2 = |z|$.

\voorbeeld{}{Bepaal de vierkantswortels van $5 + 12i$.\label{vb:h2-vw}}{%
\uitleg{We zoeken $w = x + yi$ met $w^2 = 5 + 12i$. Uit het reëel deel, het imaginair deel en de modulus halen we drie vergelijkingen voor $x$ en $y$.}
\begin{stappen}
x^2 - y^2 = 5, \quad 2xy = 12 & Reëel en imaginair deel van $w^2 = (x^2 - y^2) + 2xyi$. \\[1ex]
x^2 + y^2 = \sqrt{5^2 + 12^2} = 13 & De modulus: $x^2 + y^2 = |5 + 12i|$. \\[1ex]
2x^2 = 18, \quad 2y^2 = 8 & Tel de eerste en de derde vergelijking op, en trek ze af. \\[1ex]
x = \pm 3, \quad y = \pm 2 & Nog niet elke combinatie is goed. \\[1ex]
w = 3 + 2i \quad \text{of} \quad w = -3 - 2i & $xy = 6 > 0$: $x$ en $y$ hebben hetzelfde teken.
\end{stappen}}

Vergeet de laatste stap niet. Van de vier combinaties $\pm 3 \pm 2i$ zijn er maar twee goed: $(3 - 2i)^2 = 9 - 12i + 4i^2 = 5 - 12i$, en dat is de toegevoegde van $5 + 12i$.

\voorbeeld{}{Bepaal de vierkantswortels van $-8 - 6i$.\label{vb:h2-vw2}}{%
\uitleg{We volgen dezelfde stappen. Nu is $2xy$ negatief, en dat bepaalt welke tekens bij elkaar horen.}
\begin{stappen}
x^2 - y^2 = -8, \quad 2xy = -6, \quad x^2 + y^2 = 10 & Want $|-8 - 6i| = \sqrt{64 + 36} = 10$. \\[1ex]
x^2 = 1, \quad y^2 = 9 & Optellen en aftrekken. \\[1ex]
w = 1 - 3i \quad \text{of} \quad w = -1 + 3i & $xy = -3 < 0$: verschillend teken.
\end{stappen}
Controle: $(1 - 3i)^2 = 1 - 6i + 9i^2 = -8 - 6i$.}

\begin{werkwijze}
Stel $w = x + yi$. Schrijf op: $x^2 - y^2 = \Re z$, \ $x^2 + y^2 = |z|$ \ en \ $2xy = \Im z$. Uit de eerste twee vind je $x^2$ en $y^2$. Het teken van $xy$ zegt welke tekens van $x$ en $y$ samenhoren. Controleer door te kwadrateren.
\end{werkwijze}

\nujij{Bepaal de vierkantswortels van $3 - 4i$.}{$x^2 - y^2 = 3$, $x^2 + y^2 = 5$, $xy = -2$: $x^2 = 4$, $y^2 = 1$, dus $w = \pm(2 - i)$.}

\begin{oefeningen}
\oef{h2-vw}
\end{oefeningen}

%========================================================================================
\opdrachtenkop

\begin{oefeningen}
\oef{h2-o3}

\oef{h2-o2}

\oef{h2-o19}

\oef{h2-o4}

\oef{h2-o20}

\oef{h2-o6}

\oef{h2-o21}

\oef{h2-o5}

\oef{h2-o1}

\oef{h2-o7}

\oef{h2-o8}

\oef{h2-o22}

\oef{h2-o10}

\oef{h2-o25}

\oef{h2-o12}

\oef{h2-o9}

\oef{h2-o24}

\oef{h2-o13}

\oef{h2-o11}

\oef{h2-o23}

\oef{h2-o16}

\oef{h2-o14}

\oef{h2-o17}

\oef{h2-o26}

\oef{h2-o15}

\oef{h2-o18}
\end{oefeningen}

%----------------------------------------------------------------------------------------
% KOPPEL: verbanden met andere leerstof
%----------------------------------------------------------------------------------------
\needspace{20\baselineskip}
\opdrachtsectie{Koppel}

\lkr{Tweedegraadsfuncties kennen de leerlingen uit de 2de graad. De tweede oefening is een stukje geschiedenis: zo kwamen de complexe getallen in de wiskunde.}

\begin{oefeningen}
\oef{h2-koppel-1}

\oef{h2-koppel-2}
\end{oefeningen}

%----------------------------------------------------------------------------------------
\opdrachtsectie{Test jezelf}

\begin{oefeningen}
\oef{h2-t1}

\oef{h2-t2}

\oef{h2-t3}

\oef{h2-t4}

\oef{h2-t5}

\oef{h2-t6}

\oef{h2-t7}

\oef{h2-t8}
\end{oefeningen}

\herhaalschema{%
  \fv{oef:h2-t1}{vb:h2-wortel}{oef:h2-wortel}%
  \fv{oef:h2-t2}{vb:h2-vkv}{oef:h2-vkv}%
  \fv{oef:h2-t3}{vb:h2-vkv-a}{oef:h2-vkv}%
  \fv{oef:h2-t4}{vb:h2-opstellen}{oef:h2-opstellen}%
  \fv{oef:h2-t5}{vb:h2-ontbinden}{oef:h2-opstellen}%
  \fv{oef:h2-t6}{vb:h2-hoger}{oef:h2-hoger}%
  \fv{oef:h2-t7}{vb:h2-vw}{oef:h2-vw}%
  \fv{oef:h2-t8}{vb:h2-toegevoegd}{oef:h2-hoger3}}

\ifzelfstudie
\section*{Maak je eigen samenvatting}
Sluit dit boek en schrijf in je schrift: wat een vierkantswortel van een complex getal is, de vierkantswortels van een negatief getal, hoe je een tweedegraadsvergelijking oplost als de discriminant negatief is, de rekenregels voor de toegevoegde, waarom niet-reële oplossingen in toegevoegde paren komen, som, product en ontbinding, hoe je een vergelijking van hogere graad aanpakt, en hoe je de vierkantswortels van een complex getal vindt. Vergelijk daarna met de kaders.
\fi
